\documentclass[aspectratio=169]{beamer}

\usetheme{Madrid}
\usecolortheme{default}
\usefonttheme{professionalfonts}

\usepackage{amsmath,amssymb,bm}
\usepackage{braket}
\usepackage{tikz}
\usetikzlibrary{quantikz}

\title{Preparing Bell States from $\ket{00}$ with Quantum Circuits}
\author{}
\institute{}
\date{}

\begin{document}

\begin{frame}
  \titlepage
\end{frame}

% ==========================================================
\begin{frame}{Goal: Create Bell states starting from $\ket{00}$}
We start with the two-qubit computational basis state
\[
\ket{\psi_0}=\ket{00}.
\]

\medskip
Using a small set of gates we can prepare the four Bell states:
\[
\ket{\Phi^\pm}=\frac{\ket{00}\pm \ket{11}}{\sqrt{2}},
\qquad
\ket{\Psi^\pm}=\frac{\ket{01}\pm \ket{10}}{\sqrt{2}}.
\]

\medskip
Core idea:
\begin{itemize}
  \item Use a Hadamard on qubit 0 to create a superposition.
  \item Use an entangling gate (CNOT) to correlate qubit 1 with qubit 0.
  \item Apply simple single-qubit gates ($X$, $Z$) to obtain the other Bell states.
\end{itemize}
\end{frame}

% ==========================================================
\begin{frame}{Building blocks: single- and two-qubit gates}
\textbf{Hadamard gate} on one qubit:
\[
H\ket{0}=\frac{\ket{0}+\ket{1}}{\sqrt{2}},
\qquad
H\ket{1}=\frac{\ket{0}-\ket{1}}{\sqrt{2}}.
\]

\medskip
\textbf{CNOT} with control qubit 0 and target qubit 1:
\[
\mathrm{CNOT}\ket{c,t}=\ket{c,\,t\oplus c}.
\]
Thus, if $c=0$ nothing happens; if $c=1$, the target bit flips.

\medskip
We also use Pauli gates:
\[
X\ket{0}=\ket{1},\; X\ket{1}=\ket{0}, 
\qquad
Z\ket{0}=\ket{0},\; Z\ket{1}=-\ket{1}.
\]
\end{frame}

% ==========================================================
\begin{frame}{Circuit for $\ket{\Phi^+}$ from $\ket{00}$}
Start with $\ket{00}$.

\medskip
Step 1: Apply $H$ on qubit 0:
\[
(H\otimes I)\ket{00}=\frac{1}{\sqrt{2}}\left(\ket{00}+\ket{10}\right).
\]

\medskip
Step 2: Apply $\mathrm{CNOT}_{0\to 1}$:
\[
\mathrm{CNOT}\,\frac{\ket{00}+\ket{10}}{\sqrt{2}}
=\frac{1}{\sqrt{2}}\left(\ket{00}+\ket{11}\right)=\ket{\Phi^+}.
\]

\medskip
\begin{center}
\begin{quantikz}
\lstick{$\ket{0}$} & \gate{H} & \ctrl{1} & \qw \\
\lstick{$\ket{0}$} & \qw      & \targ{}  & \qw
\end{quantikz}
\end{center}
\end{frame}

% ==========================================================
\begin{frame}{Circuit for $\ket{\Phi^-}$ from $\ket{00}$}
Use the same entangling circuit as for $\ket{\Phi^+}$, then add a $Z$ phase flip.

\medskip
Since
\[
(I\otimes Z)\ket{\Phi^+}=\frac{1}{\sqrt{2}}\left(\ket{00}-\ket{11}\right)=\ket{\Phi^-},
\]
we can implement $\ket{\Phi^-}$ by applying $Z$ to either qubit (global phase conventions aside).

\begin{center}
\begin{quantikz}
\lstick{$\ket{0}$} & \gate{H} & \ctrl{1} & \gate{Z} & \qw \\
\lstick{$\ket{0}$} & \qw      & \targ{}  & \qw      & \qw
\end{quantikz}
\end{center}

\medskip
Equivalent variant (apply $Z$ on the second qubit):
\[
(I\otimes Z)\ket{\Phi^+}=\ket{\Phi^-}.
\]
\end{frame}

% ==========================================================
\begin{frame}{Circuit for $\ket{\Psi^+}$ from $\ket{00}$}
One convenient route:
\[
\ket{\Psi^+} = (I\otimes X)\ket{\Phi^+}.
\]

\medskip
Thus: build $\ket{\Phi^+}$ and then flip the second qubit with $X$:
\[
(I\otimes X)\frac{\ket{00}+\ket{11}}{\sqrt{2}}
=\frac{\ket{01}+\ket{10}}{\sqrt{2}}=\ket{\Psi^+}.
\]

\begin{center}
\begin{quantikz}
\lstick{$\ket{0}$} & \gate{H} & \ctrl{1} & \qw      & \qw \\
\lstick{$\ket{0}$} & \qw      & \targ{}  & \gate{X} & \qw
\end{quantikz}
\end{center}

\medskip
(You can also apply $X$ before the CNOT on the target qubit; the above is the clearest construction.)
\end{frame}

% ==========================================================
\begin{frame}{Circuit for $\ket{\Psi^-}$ from $\ket{00}$}
A simple construction uses
\[
\ket{\Psi^-}=(I\otimes X)\ket{\Phi^-}
\quad\text{or}\quad
\ket{\Psi^-}=(I\otimes Z)(I\otimes X)\ket{\Phi^+}.
\]

\medskip
Starting from $\ket{\Phi^+}$:
\[
\ket{\Phi^+}\xrightarrow{\,X\text{ on qubit 1}\,}\ket{\Psi^+}
\xrightarrow{\,Z\text{ on qubit 0 (or 1)}\,}\ket{\Psi^-}.
\]

\begin{center}
\begin{quantikz}
\lstick{$\ket{0}$} & \gate{H} & \ctrl{1} & \gate{Z} & \qw \\
\lstick{$\ket{0}$} & \qw      & \targ{}  & \gate{X} & \qw
\end{quantikz}
\end{center}

\medskip
Up to a global phase, applying $Z$ on either qubit produces the relative minus sign between $\ket{01}$ and $\ket{10}$.
\end{frame}

% ==========================================================
\begin{frame}{Summary: One entangling core circuit + local Pauli gates}
Core entangling circuit prepares $\ket{\Phi^+}$:
\[
\ket{00}\;\xrightarrow{H\otimes I}\;\frac{\ket{00}+\ket{10}}{\sqrt{2}}
\;\xrightarrow{\mathrm{CNOT}_{0\to 1}}\;\ket{\Phi^+}.
\]

\medskip
Then obtain the other Bell states by local operations:
\[
\ket{\Phi^-}=(Z\otimes I)\ket{\Phi^+},
\qquad
\ket{\Psi^+}=(I\otimes X)\ket{\Phi^+},
\qquad
\ket{\Psi^-}=(Z\otimes I)(I\otimes X)\ket{\Phi^+}.
\]

\medskip
\begin{center}
\begin{quantikz}
\lstick{$\ket{0}$} & \gate{H} & \ctrl{1} & \qw \\
\lstick{$\ket{0}$} & \qw      & \targ{}  & \qw
\end{quantikz}
\end{center}
\end{frame}

% ==========================================================
\begin{frame}{Optional: Verify with state evolution (conceptual)}
If we denote the circuit unitary by $U$, then
\[
U\ket{00}=\ket{\Phi^+}.
\]
Local Pauli operations generate the full Bell basis:
\[
\left\{\ket{\Phi^+},\; (Z\otimes I)\ket{\Phi^+},\; (I\otimes X)\ket{\Phi^+},\; (Z\otimes I)(I\otimes X)\ket{\Phi^+}\right\}.
\]

\medskip
In quantum information, this is summarized as:
\[
\text{Bell basis}=\{(I\otimes \sigma)\ket{\Phi^+}\;:\;\sigma\in\{I,X,Z,XZ\}\}.
\]
\end{frame}

\end{document}
